\documentclass[12pt]{article}

\usepackage{amsmath}
\usepackage{amssymb}
\usepackage{amsthm}
\usepackage{mathtools}
\usepackage{mathrsfs}
\usepackage{hyperref}
\usepackage{braket}
\usepackage{bm}
\usepackage{graphicx}
\usepackage[
    left = 1in,
    right = 1in,
    top = 1in,
    bottom = 1in
]{geometry}
\usepackage{titling}

\setlength{\droptitle}{-0.7in}

\begin{document}

\title{Documentation}
\author{}
\date{}
\maketitle
\vspace{-1in}

\section*{Introduction \& Topic}

For this project, our group opted to construct a \textbf{cone} with 
maximum radius $32"$ and height, from the xy-plane, $36"$ (or $2.5"$ 
and $3"$, respectively). The cone was limited to only positive $z$
values, for the sake of construction, and no other domain restrictions
were imposed.

The official equation used was

\[
\frac{x^{2}}{900} + \frac{y^{2}}{900} = \frac{z^{2}}{1296}, \ z > 0,
\]

and we …

\section*{Construction Process}

…

\newpage

\title{Calculations}
\author{}
\date{}
\maketitle
\vspace{-1in}

\section*{Precise Theoretic Value}

To calculate the precise theoretic value, we will use an integral,
specifically a double integral.

\section*{Approximate Theoretic Value}

To calculate the approximate theoretic value, we will use a Riemann 
summation, specifically two summations.

\newpage

\title{Discussion}
\author{}
\date{}
\maketitle
\vspace{-1in}

\section*{Precise Theoretic Value}

…

\section*{Approximate Theoretic Value}

…

\end{document}
